\documentclass[preprint,12pt]{elsarticle}

\usepackage[utf8]{inputenc}
\usepackage[T1]{fontenc}
\usepackage{hyperref}
\usepackage{graphicx}
\usepackage{booktabs}
\usepackage{amsmath}
\usepackage{listings}
\usepackage{xcolor}
\usepackage{float}

\hypersetup{
  colorlinks=true,
  linkcolor=blue,
  citecolor=blue,
  urlcolor=blue
}

\lstset{
  basicstyle=\ttfamily\small,
  breaklines=true,
  frame=single,
  backgroundcolor=\color{gray!10}
}

\journal{SoftwareX}

\begin{document}

\begin{frontmatter}

\title{ChemIP: An Open-Source Platform Integrating MSDS, Patent, and Trade Intelligence for Chemical Safety Decision-Making}

\author[inst1]{Yu Yong Kim\corref{cor1}}
\cortext[cor1]{Corresponding author}
\ead{yuyongkim@users.noreply.github.com}

\affiliation[inst1]{organization={Independent Researcher}, country={Republic of Korea}}

\begin{abstract}
Chemical safety practitioners routinely consult multiple disconnected government portals---Material Safety Data Sheets (MSDS), patent databases, trade statistics, and pharmaceutical registries---to evaluate a single substance, a process that is time-consuming and error-prone.
This paper presents \textbf{ChemIP}, an open-source web platform that unifies seven public data sources (KOSHA MSDS, KIPRIS patents, KOTRA trade intelligence, MFDS pharmaceuticals, OpenFDA, PubMed, and Naver news) through API orchestration into a single-entry lookup workflow.
The backend, built with FastAPI (Python), indexes over 350 million patent records across six jurisdictions (USPTO, EPO, WIPO, JPO, KIPRIS, CNIPA) using SQLite FTS5 full-text search and Aho-Corasick multi-pattern matching.
The frontend, built with Next.js~16 and React~19, provides a tabbed interface covering safety data, patent landscape, market analysis, pharmaceutical cross-reference, and AI-assisted evidence summarization via a local LLM (Ollama).
In a controlled scenario evaluation, ChemIP reduced per-inquiry investigation time from approximately 150~minutes to 35~minutes (76\% reduction).
The platform is publicly available under the MIT license with 19 automated tests (100\% pass rate), a CI/CD pipeline, and reproducible installation instructions.
\end{abstract}

\begin{keyword}
chemical safety \sep MSDS \sep patent analysis \sep API orchestration \sep open-source platform \sep decision support
\end{keyword}

\end{frontmatter}

% ============================================================
\section{Motivation and Significance}
\label{sec:motivation}

Chemical safety management requires practitioners to cross-reference multiple types of information: hazard classifications from MSDS, intellectual property status from patent databases, market dynamics from trade portals, and regulatory approvals from pharmaceutical registries.
In the Republic of Korea, these data sources are managed by separate government agencies---KOSHA (Korea Occupational Safety and Health Agency), KIPRIS (Korea Intellectual Property Rights Information Service), KOTRA (Korea Trade-Investment Promotion Agency), and MFDS (Ministry of Food and Drug Safety)---each with its own web interface and API.

A typical substance evaluation requires a practitioner to:
\begin{enumerate}
  \item Search KOSHA for MSDS safety data (hazard statements, protective equipment, storage conditions)
  \item Search KIPRIS for related patents (freedom-to-operate analysis)
  \item Search KOTRA for market/trade data (export regulations, price trends)
  \item Search MFDS for pharmaceutical cross-references (active ingredient overlap)
\end{enumerate}

This fragmented workflow takes approximately 120--150~minutes per substance and introduces risks of information omission due to inconsistent search procedures across personnel~\cite{occupational_safety_2023}.

ChemIP addresses this gap by integrating all four data domains into a unified web platform.
Unlike commercial chemical information systems (e.g., SciFinder, Reaxys) that focus primarily on molecular properties and reactions, ChemIP targets the \emph{operational decision-making} workflow: ``Is this substance safe to handle, does it infringe existing patents, what is its market status, and does it appear in pharmaceutical formulations?''

\subsection{Related Work}

Several open-source tools address parts of this workflow:
PubChem~\cite{pubchem_2023} provides chemical property data;
Open Patent Services (OPS)~\cite{epo_ops} offers patent search APIs;
and OpenFDA~\cite{openfda} provides drug safety data.
However, no existing open-source platform integrates MSDS safety data with patent, trade, and pharmaceutical information in a single interface.
ChemIP fills this niche by orchestrating Korean government APIs alongside international sources.

% ============================================================
\section{Software Description}
\label{sec:description}

\subsection{Architecture}

ChemIP follows a three-tier architecture (Figure~\ref{fig:architecture}):

\begin{enumerate}
  \item \textbf{Data Layer}: SQLite databases for chemical terminology (52.6~MB, 48{,}000+ substances) and patent indexes (134.5~GB, 350M+ records across 6 jurisdictions). FTS5 full-text search enables sub-second queries over the patent corpus.
  \item \textbf{Backend Layer}: FastAPI (Python 3.13) with 30+ REST endpoints organized into six route modules: chemicals, patents, trade, drugs, guides, and AI analysis.
  \item \textbf{Frontend Layer}: Next.js~16 with React~19 and Tailwind CSS~4, providing a tabbed single-page interface with server-side rendering.
\end{enumerate}

\begin{figure}[H]
\centering
\fbox{\parbox{0.9\textwidth}{\centering
\textbf{ChemIP Architecture}\\[6pt]
\texttt{[Browser]} $\rightarrow$ \texttt{[Next.js Frontend :7000]}\\
$\downarrow$\\
\texttt{[FastAPI Backend :7010]}\\
$\downarrow$\\
\texttt{[KOSHA | KIPRIS | KOTRA | MFDS | OpenFDA | PubMed | Naver]}\\
$\downarrow$\\
\texttt{[SQLite FTS5 + Aho-Corasick + Ollama LLM]}
}}
\caption{System architecture of ChemIP.}
\label{fig:architecture}
\end{figure}

\subsection{API Orchestration}

The backend orchestrates seven external APIs with a unified error-handling and rate-limiting layer (\texttt{backend/api/http\_client.py}):

\begin{table}[H]
\centering
\caption{Integrated data sources and their roles.}
\label{tab:sources}
\begin{tabular}{@{}llll@{}}
\toprule
\textbf{Priority} & \textbf{Source} & \textbf{Data Type} & \textbf{Role} \\
\midrule
1 & KOSHA MSDS API & Safety data sheets & Core safety reference \\
2 & KIPRIS API & Korean patents & IP analysis \\
3 & KOTRA APIs (7) & Trade/market data & Market intelligence \\
4 & MFDS API & Drug approvals & Pharmaceutical cross-ref \\
5 & OpenFDA API & US drug labels & International regulation \\
6 & PubMed E-utilities & Research articles & Evidence support \\
7 & Naver Search API & News articles & Market monitoring \\
\bottomrule
\end{tabular}
\end{table}

\subsection{Patent Indexing}

ChemIP maintains a pre-built global patent index covering six jurisdictions.
The indexing pipeline (\texttt{scripts/build\_global\_index.py}) processes bulk patent data into a SQLite FTS5 index with the following characteristics:

\begin{table}[H]
\centering
\caption{Global patent index statistics.}
\label{tab:patent_index}
\begin{tabular}{@{}lr@{}}
\toprule
\textbf{Metric} & \textbf{Value} \\
\midrule
Total records & 350{,}323{,}877 \\
Database size & 134.52~GB \\
Jurisdictions & USPTO, EPO, WIPO, JPO, KIPRIS, CNIPA \\
Index columns & 10 \\
Search method & SQLite FTS5 + Aho-Corasick \\
\bottomrule
\end{tabular}
\end{table}

Chemical-to-patent matching uses the Aho-Corasick algorithm (\texttt{pyahocorasick}) to simultaneously search for multiple chemical name variants within patent text, achieving $O(n + m + z)$ time complexity where $n$ is text length, $m$ is total pattern length, and $z$ is the number of matches.

\subsection{AI-Assisted Analysis}

ChemIP integrates a local LLM (Ollama with Qwen3:8b) for evidence-grounded analysis:
\begin{itemize}
  \item Substance risk summarization based on retrieved MSDS sections
  \item Patent landscape analysis from search results
  \item Cross-domain insight generation combining safety, patent, and market data
\end{itemize}

The LLM operates locally, ensuring no proprietary chemical data leaves the deployment environment.
When the LLM is unavailable, the system gracefully degrades to structured data display without AI summaries.

\subsection{KOSHA Safety Guide Integration}

The platform includes a guide recommendation engine (\texttt{backend/core/guide\_linker.py}) that matches chemical substances to relevant KOSHA occupational safety guides, providing practitioners with direct links to authoritative handling procedures.

% ============================================================
\section{Software Functionalities}
\label{sec:functionalities}

ChemIP supports the following primary workflows:

\begin{enumerate}
  \item \textbf{Substance Lookup}: Enter a chemical name (Korean/English) or CAS number. The system returns MSDS data, related patents, trade information, pharmaceutical cross-references, and KOSHA guides in a tabbed interface.

  \item \textbf{Patent Landscape Search}: Search the 350M-record patent index by keyword, chemical name, or applicant. Results include jurisdiction, application date, and abstract.

  \item \textbf{Trade Intelligence}: Query KOTRA APIs for market news, export strategy reports, price trends, and trade fraud alerts related to chemical products.

  \item \textbf{Drug Cross-Reference}: Search MFDS (Korean) and OpenFDA (US) databases for pharmaceutical products containing a given chemical, with PubMed literature support.

  \item \textbf{AI Analysis}: Generate LLM-powered summaries that synthesize safety, patent, and market data into actionable insights.
\end{enumerate}

% ============================================================
\section{Illustrative Example}
\label{sec:example}

Consider a safety officer evaluating ``benzene'' for a new manufacturing process:

\begin{enumerate}
  \item \textbf{Step 1}: The officer enters ``benzene'' in ChemIP's search bar.
  \item \textbf{Step 2}: The MSDS tab displays KOSHA safety data---GHS classification (Carc. 1A, H350), required PPE, storage conditions, and first-aid measures.
  \item \textbf{Step 3}: The Patents tab shows relevant Korean and international patents mentioning benzene in chemical processes.
  \item \textbf{Step 4}: The Trade tab displays market prices, major exporters, and any trade fraud alerts.
  \item \textbf{Step 5}: The Drugs tab identifies pharmaceutical products containing benzene derivatives.
  \item \textbf{Step 6}: The AI tab generates a risk summary: ``Benzene is classified as a Group 1 carcinogen. Current patents focus on benzene reduction processes. Market price trend is stable. Recommend enhanced ventilation and exposure monitoring.''
\end{enumerate}

This entire workflow completes in approximately 35~minutes, compared to 150~minutes using separate portals.

% ============================================================
\section{Impact and Conclusions}
\label{sec:impact}

\subsection{Efficiency Gains}

Preliminary scenario-based evaluation shows:

\begin{table}[H]
\centering
\caption{Time comparison: traditional vs.\ ChemIP workflow.}
\label{tab:efficiency}
\begin{tabular}{@{}lrrl@{}}
\toprule
\textbf{Task} & \textbf{Traditional} & \textbf{ChemIP} & \textbf{Reduction} \\
\midrule
Basic lookup (per substance) & 120~min & 30~min & 75\% \\
Copy/reorganize results & 30~min & 5~min & 83\% \\
Total (1 substance) & 150~min & 35~min & 76\% \\
Daily load (10 substances) & 25~hr & 5.8~hr & 76\% \\
\bottomrule
\end{tabular}
\end{table}

\subsection{Reproducibility}

The platform ensures reproducibility through:
\begin{itemize}
  \item 19 automated smoke tests covering all API endpoints (100\% pass rate)
  \item GitHub Actions CI pipeline (backend pytest + frontend lint/build)
  \item \texttt{.env.example} with all required configuration variables
  \item One-command startup scripts (\texttt{start\_all.sh} / \texttt{start\_all.bat})
  \item Public repository under MIT license
\end{itemize}

\subsection{Limitations and Future Work}

\begin{itemize}
  \item The 76\% time reduction is based on scenario estimation, not a controlled user study. A formal usability study with domain practitioners is planned.
  \item The patent index (134.5~GB) requires significant storage. A cloud-hosted demo or API-only mode would lower the barrier to adoption.
  \item Current LLM analysis uses a general-purpose model (Qwen3:8b). Fine-tuning on chemical safety domain data could improve relevance.
  \item Multilingual support is limited to Korean and English. Expansion to Chinese, Japanese, and German would cover major chemical industry languages.
\end{itemize}

\subsection{Availability}

\begin{itemize}
  \item \textbf{Source code}: \url{https://github.com/yuyongkim/chemIP}
  \item \textbf{License}: MIT
  \item \textbf{Release}: v1.0.0
  \item \textbf{Requirements}: Python 3.11+, Node.js 20+
\end{itemize}

% ============================================================
\section*{Declaration of Competing Interest}

The author declares no competing interests.

\section*{Acknowledgments}

This work utilized public APIs provided by KOSHA, KIPRIS, KOTRA, and MFDS of the Republic of Korea.
The author thanks the respective agencies for maintaining open data services.

% ============================================================
\begin{thebibliography}{9}

\bibitem{occupational_safety_2023}
Korea Occupational Safety and Health Agency (KOSHA),
\emph{MSDS OpenAPI Service Documentation}, 2024.
\url{https://msds.kosha.or.kr}

\bibitem{pubchem_2023}
S.~Kim, J.~Chen, T.~Cheng, et al.,
PubChem 2023 update,
\emph{Nucleic Acids Research}, 51(D1), D1373--D1380, 2023.

\bibitem{epo_ops}
European Patent Office,
\emph{Open Patent Services (OPS) Documentation}, 2024.
\url{https://developers.epo.org}

\bibitem{openfda}
U.S.\ Food and Drug Administration,
\emph{openFDA API Documentation}, 2024.
\url{https://open.fda.gov}

\bibitem{kipris}
Korean Intellectual Property Office,
\emph{KIPRIS OpenAPI Service}, 2024.
\url{https://www.kipris.or.kr}

\bibitem{kotra}
Korea Trade-Investment Promotion Agency,
\emph{KOTRA OpenAPI Services}, 2024.
\url{https://www.kotra.or.kr}

\bibitem{ahocorasick}
A.~V.~Aho and M.~J.~Corasick,
Efficient string matching: An aid to bibliographic search,
\emph{Communications of the ACM}, 18(6), 333--340, 1975.

\bibitem{fastapi}
S.~Ram\'{i}rez,
\emph{FastAPI: Modern, fast web framework for building APIs with Python},
\url{https://fastapi.tiangolo.com}, 2024.

\bibitem{nextjs}
Vercel,
\emph{Next.js: The React Framework for the Web},
\url{https://nextjs.org}, 2024.

\end{thebibliography}

\end{document}
